\documentclass{article}
\usepackage{graphicx} % Required for inserting images

\title{DDMS}
\author{אורן רם}
\date{July 2026}

\begin{document}





\section{The Likelihood Ratio as a Sufficient Statistic}

\subsection{Definition}
For each sensor $i$, define the likelihood ratio of its observation as
\begin{equation}
    l^i := L(y^i) := \frac{f_{Y^i|H}(y^i \mid H_2)}{f_{Y^i|H}(y^i \mid H_1)}
    \quad \text{a.s.}
    \label{eq:likelihood-ratio}
\end{equation}
a single real number capturing how much more consistent the raw observation $y^i$
is with $H_2$ than with $H_1$.

\subsection{Why $l^i$ is sufficient}
$l^i$ is a sufficient statistic for $H$: once its value is known, nothing else
about the raw observation $y^i$ can further move one's belief about which
hypothesis is true. With only two hypotheses on the table, the entire way in which
the observation's density can vary with $H$ is exhausted by a single comparison -
how large the density is under $H_2$ relative to how large it is under $H_1$, at
that particular observation. A single number is therefore inherently enough to
hold everything that changes with $H$, and that number is, by construction,
exactly $l^i$. (How this shows up concretely inside the paper's own machinery -
via Lemma~1 -- is worth treating separately once we get there.)

\subsection{No information loss: a Markov-chain and DPI reading}
Sufficiency means that any two raw observations producing the same $l^i$ are, for
the purpose of distinguishing $H_1$ from $H_2$, interchangeable: whatever detail of
$y^i$ is discarded in forming $l^i$ carries no evidential weight about $H$. This
has a clean information-theoretic statement. Since $l^i$ is a deterministic
function of $y^i$, the chain
\[
    H \;\to\; y^i \;\to\; l^i
\]
is trivially Markov, and the Data Processing Inequality (DPI) gives
\begin{equation}
    I(H; l^i) \le I(H; y^i).
\end{equation}
Equality holds \emph{iff} $l^i$ is sufficient for $H$ -- equivalently, iff
$H \to l^i \to y^i$ is \emph{also} a Markov chain, i.e.\ $y^i$ carries no further
information about $H$ once $l^i$ is fixed, gives equality
\begin{equation}
    I(H; l^i) = I(H; y^i).
\end{equation}
\end{document}
